\cleardoublepage % memastikan bab baru dimulai di halaman ganjil (kanan)
\chapter{RENCANA SELANJUTNYA}
\label{chap:rencana-selanjutnya}

Bab ini menjabarkan cara solusi pada Bab~IV akan dibangun dan diuji. Pembahasan mencakup rencana implementasi, desain pengujian dan evaluasi, serta analisis risiko dan mitigasinya. Setiap rencana dikaitkan dengan kebutuhan pada Subbab~\ref{sec:analisis_kebutuhan} dan kriteria keberhasilan pada Subbab~\ref{sec:tujuan}.

\section{Rencana Implementasi}
\label{sec:rencana_implementasi}

\subsection{Tahapan dan Jadwal}
\label{subsec:jadwal}
Implementasi direncanakan selama 16 minggu pada semester kedua tugas akhir. Tahapan dan jadwalnya disajikan pada Tabel~\ref{tab:jadwal}. Ada tiga gerbang keputusan. Gerbang G0 berada setelah simulator selesai dan sebelum pelatihan agen. Gerbang ini mengukur kemampuan \textit{MacBook} dan menentukan kegiatan mana yang tetap di \textit{MacBook} dan mana yang dipindah ke infrastruktur eksternal (Subbab~\ref{subsec:infrastruktur}). Gerbang G1 berada setelah validasi simulator, dan pelatihan agen baru dimulai jika simulator lolos KNF-01. Gerbang G2 berada setelah pelatihan ko-evolusioner, dan eksperimen utama baru dimulai jika pelatihan memenuhi kriteria konvergensi pada Subbab~\ref{subsec:konfigurasi}. Gerbang G1 dan G2 mencegah hasil eksperimen dibangun di atas simulator atau kebijakan yang belum layak.

\begin{table}[htbp]
\centering
\caption{Tahapan dan jadwal implementasi}
\label{tab:jadwal}
\small
\begin{tabular}{|p{0.7cm}|p{3.0cm}|p{4.5cm}|p{1.4cm}|p{2.2cm}|}
\hline
\textbf{No.} & \textbf{Tahapan} & \textbf{Kegiatan utama} & \textbf{Minggu} & \textbf{Kebutuhan} \\
\hline
1 & Perancangan rinci & Diagram C4 dan UML, spesifikasi \textit{state}, \textit{action}, dan \textit{reward} final & 1--2 & KF-04 \\
\hline
2 & Implementasi simulator & Modul validator, jaringan, pengelola suara, pemilihan \textit{proposer}, dan insentif, beserta uji unit & 2--5 & KF-01 sampai KF-03 \\
\hline
3 & Verifikasi, kalibrasi, dan benchmark & Reproduksi hasil \textcite{laneve2025stochastic}, pengambilan data \textit{on-chain}, uji Kolmogorov-Smirnov (gerbang G1), uji asap pelatihan dan pengukuran kecepatan \textit{MacBook} Reproduksi hasil \textcite{laneve2025stochastic}, pengambilan data \textit{on-chain}, uji Kolmogorov-Smirnov (gerbang G1) & 5--6 & KNF-01 \\
\hline
4 & Pembanding & Pertahanan statis, berbasis aturan, dan DRL satu sisi & 7--8 & KF-07 \\
\hline
5 & Pelatihan ko-evolusioner & Pelatihan kedua agen, penyetelan bobot \textit{reward}, pemeriksaan konvergensi (gerbang G2) & 8--11 & KF-05, KF-06 \\
\hline
6 & Eksperimen dan evaluasi & Eksperimen utama, uji sensitivitas, ablasi, uji skala & 11--13 & KF-07, KF-08 \\
\hline
7 & Analisis dan penulisan & Uji hipotesis, analisis strategi, penyusunan laporan, persiapan sidang & 13--16 & Seluruh RQ \\
\hline
\end{tabular}
\end{table}

\subsection{Alat dan Bahan}
\label{subsec:alat_bahan}
Alat dan bahan yang diperlukan disajikan pada Tabel~\ref{tab:alat_bahan}. Pemilihan perangkat lunak mengikuti lingkungan yang terbukti dapat mereproduksi pelatihan DRL pada \textcite{skowronski2026adaptive}, yaitu Python dan PyTorch. PettingZoo ditambahkan karena arena melibatkan dua agen. Perangkat utama adalah \textit{MacBook Pro} dengan Apple M5 Pro milik penulis. Perangkat eksternal hanya disiapkan sebagai cadangan, dengan acuan kebutuhan pada \textcite{skowronski2026adaptive} untuk pelatihan dan \textcite{laneve2025stochastic} untuk simulasi berulang. Pembagian tugas keduanya dijelaskan pada Subbab~\ref{subsec:infrastruktur}.

\begin{table}[htbp]
\centering
\caption{Alat dan bahan implementasi}
\label{tab:alat_bahan}
\small
\begin{tabular}{|p{2.3cm}|p{4.4cm}|p{5.5cm}|}
\hline
\textbf{Kategori} & \textbf{Alat atau bahan} & \textbf{Fungsi} \\
\hline
Bahasa dan pustaka utama & Python 3.10 atau lebih baru, PyTorch 2.x (backend MPS atau CPU di \textit{MacBook}, CUDA di infrastruktur eksternal) & Implementasi simulator dan jaringan DDQN \\
\hline
Antarmuka multi-agen & PettingZoo & Antarmuka \textit{Markov game} dua agen dengan frekuensi keputusan berbeda \\
\hline
Komputasi ilmiah & NumPy (termasuk \textit{SeedSequence}), SciPy, Pandas & Pembangkitan \textit{seed}, uji Kolmogorov-Smirnov, uji Wilcoxon, olah data \\
\hline
Pemantauan dan versi & TensorBoard, Git, GitHub & Pemantauan kurva pelatihan, pengelolaan versi kode dan konfigurasi \\
\hline
Data & API \textit{beaconcha.in}, \textit{BigQuery}, \textit{consensus-specs} & Statistik \textit{slot} kosong, partisipasi \textit{attestation}, blok yatim, dan nilai parameter acuan \\
\hline
Perangkat utama & \textit{MacBook Pro} Apple M5 Pro, RAM terpadu dengan kapasitas dicatat pada gerbang G0 & Pengembangan, simulasi paralel multi-inti, verifikasi, validasi, evaluasi, dan pelatihan sejauh lolos gerbang G0 \\
\hline
Perangkat eksternal (cadangan) & GPU NVIDIA 24 GB (fasilitas kampus atau sewa awan) dengan CUDA & Hanya untuk kegiatan yang tidak lolos gerbang G0. Acuan: satu GPU RTX 3090 \parencite{skowronski2026adaptive} dan mesin 8 vCPU 64 GB \parencite{laneve2025stochastic} \\
\hline
Penyimpanan cadangan & Hardisk eksternal atau penyimpanan awan & Salinan konfigurasi, \textit{checkpoint}, dan log pelatihan \\
\hline
\end{tabular}
\end{table}
\subsection{Pembagian Infrastruktur}
\label{subsec:infrastruktur}
\textit{MacBook Pro} M5 Pro digunakan lebih dulu untuk seluruh pekerjaan. Infrastruktur eksternal baru dipakai jika \textit{MacBook} terbukti tidak mampu menyelesaikan suatu kegiatan sesuai jadwal. Kemampuan itu diukur pada gerbang G0, bukan diperkirakan. Pembagian ini didasarkan pada empat pertimbangan.
\begin{enumerate}
	\item Simulator tidak memerlukan GPU. Simulator Gasper berjalan di CPU, dan \textcite{laneve2025stochastic} menjalankannya pada mesin 8 vCPU. Inti CPU \textit{MacBook} dapat dipakai paralel, sehingga tidak ada alasan teknis memindahkan simulator.
	\item Jaringan agen tergolong kecil. Arsitektur \textit{dueling} [128, 128, 64] ringan dibanding model yang lazim memerlukan GPU besar, dan satu sesi pelatihan pada \textcite{skowronski2026adaptive} sekitar 8 jam pada satu GPU. Karena itu, GPU Apple melalui backend MPS (\textit{Metal Performance Shaders}) atau CPU berpeluang memadai. Kepastiannya baru diperoleh dari pengukuran.
	\item \textit{MacBook} memiliki keterbatasan. Perangkat ini tidak mendukung CUDA, sehingga operator PyTorch tertentu mungkin belum tersedia pada MPS. Sebagai laptop, pelatihan berhari-hari juga rentan terhadap penurunan kecepatan akibat panas, dan perangkatnya hanya satu.
	\item Perbandingan harus adil. Hasil pelatihan pada MPS dan CUDA dapat berbeda sedikit karena implementasi numeriknya berbeda. Kegiatan yang hasilnya dibandingkan berpasangan pada Subbab~\ref{subsec:analisis_statistik} karena itu harus dilatih pada jenis perangkat yang sama.
\end{enumerate}

Pembagian kegiatan disajikan pada Tabel~\ref{tab:infrastruktur}. Angka jam pada kelompok B mengacu pada Tabel~\ref{tab:biaya}.

\begin{table}[htbp]
\centering
\caption{Pembagian kegiatan antara \textit{MacBook} dan infrastruktur eksternal}
\label{tab:infrastruktur}
\small
\begin{tabular}{|p{2.0cm}|p{4.4cm}|p{2.4cm}|p{3.0cm}|}
\hline
\textbf{Kelompok} & \textbf{Kegiatan} & \textbf{Infrastruktur} & \textbf{Alasan} \\
\hline
A. Pasti di \textit{MacBook} & Perancangan, implementasi simulator, uji unit VF-01 sampai VF-04, verifikasi VF-05, validasi Kolmogorov-Smirnov, pembanding statis dan berbasis aturan, pengambilan data \textit{on-chain}, evaluasi kebijakan beku, analisis statistik, penulisan & \textit{MacBook} & Hanya memerlukan CPU. Evaluasi di CPU membuat pengukuran akhir sama apa pun perangkat pelatihannya \\
\hline
B. \textit{MacBook} jika lolos G0 & Pembanding DRL satu sisi (80 jam), pelatihan ko-evolusioner utama (240 jam), penyerang yang dilatih khusus (80 jam), ablasi (120 jam) & \textit{MacBook}, dipindah jika gagal G0 & Memerlukan pelatihan jaringan, dan kecepatan MPS atau CPU belum diketahui sebelum diukur \\
\hline
C. Eksternal, bersyarat & Kegiatan kelompok B yang gagal G0, dan uji skala 1.000 validator jika RAM tidak cukup & GPU NVIDIA kampus atau awan, atau mesin RAM besar & Hanya cadangan, karena menambah biaya dan waktu penyiapan \\
\hline
\end{tabular}
\end{table}

Keputusan pada gerbang G0 mengikuti aturan berikut.
\begin{enumerate}
	\item Ukur kecepatan pelatihan (langkah per detik) pada \textit{MacBook} dengan backend MPS dan CPU, lalu pilih yang tercepat. Ukur juga kebutuhan memori simulator pada 128 dan 256 validator, lalu ekstrapolasikan ke 1.000 validator.
	\item Proyeksikan jam setiap kegiatan kelompok B dari kecepatan terukur dan jumlah langkah hingga konvergen pada uji asap.
	\item Suatu kegiatan tetap di \textit{MacBook} jika proyeksi jamnya, dibagi 0,7 untuk memperhitungkan jeda akibat panas dan pemakaian lain, tidak melebihi alokasi kalendernya pada Tabel~\ref{tab:jadwal}. Jika melebihi, kegiatan dipindah ke infrastruktur eksternal.
	\item Uji skala 1.000 validator tetap di \textit{MacBook} jika proyeksi memorinya paling banyak 75\% dari RAM. Jika lebih, uji ini dipindah ke mesin dengan RAM lebih besar.
\end{enumerate}

Dua aturan menjaga agar hasil tetap sebanding. Pertama, pembanding DRL satu sisi, pelatihan ko-evolusioner utama, dan ablasi berada pada jenis perangkat yang sama karena hasilnya dibandingkan berpasangan. Jika salah satunya gagal G0, ketiganya dipindah bersama atau cakupan kegiatan yang gagal dikurangi sesuai R8. Kedua, evaluasi kebijakan beku selalu dijalankan di CPU \textit{MacBook} dengan presisi \textit{float32}. Spesifikasi perangkat, versi Python dan PyTorch, backend, dan \textit{master seed} dicatat untuk setiap eksperimen.

\subsection{Lingkungan dan Konfigurasi}
\label{subsec:konfigurasi}
Konfigurasi awal disajikan pada Tabel~\ref{tab:konfigurasi}. Nilai hiperparameter agen diambil dari konfigurasi yang sudah teruji pada \textcite{skowronski2026adaptive}, sedangkan parameter kumpulan penyerang diambil dari \textcite{wushishi2026coadam}. Nilai ini menjadi titik awal dan dapat disetel ulang jika pelatihan tidak konvergen. Rentang parameter konsensus pada Tabel~\ref{tab:batas_parameter} dikalibrasi pada tahap ini dengan syarat perilaku normal simulator tetap lolos KNF-01.

\begin{table}[htbp]
\centering
\caption{Konfigurasi awal arena}
\label{tab:konfigurasi}
\small
\begin{tabular}{|p{4.6cm}|p{4.2cm}|p{3.4cm}|}
\hline
\textbf{Komponen} & \textbf{Nilai awal} & \textbf{Dasar} \\
\hline
Arsitektur jaringan kedua agen & \textit{Dueling}, lapisan [128, 128, 64], ReLU & \textcite{skowronski2026adaptive} \\
\hline
Laju belajar dan pengoptimal & $10^{-4}$, Adam & \textcite{skowronski2026adaptive} \\
\hline
Faktor diskon $\gamma$ & 0,99 & \textcite{skowronski2026adaptive} \\
\hline
Ukuran \textit{batch} dan \textit{replay buffer} & 64 dan 50.000 transisi & \textcite{skowronski2026adaptive} \\
\hline
Pembaruan jaringan target & Setiap 1.000 langkah & \textcite{skowronski2026adaptive} \\
\hline
Eksplorasi $\epsilon$ & 1,0 menuju 0,1 dalam 100.000 langkah & \textcite{skowronski2026adaptive} \\
\hline
PER & Eksponen prioritas 0,6, koreksi bobot 0,4 menuju 1,0 & \textcite{skowronski2026adaptive} \\
\hline
Kumpulan penyerang & $K = 10$, $p_{\text{pool}} = 0{,}30$, potret setiap 500 episode & \textcite{wushishi2026coadam} \\
\hline
Panjang episode & 100 \textit{epoch} (3.200 \textit{slot}) & Mengikuti 100 langkah per episode \parencite{wushishi2026coadam} \\
\hline
Jumlah validator & 13 (verifikasi), 64, 128, 256 (utama), 1.000 (uji skala) & Subbab~\ref{sec:ruang_lingkup} \\
\hline
Porsi \textit{stake} penyerang & 0,10; 0,20; 0,30 & Di bawah sepertiga (Subbab~\ref{subsec:ambang}) \\
\hline
\textit{Seed} & Satu \textit{master seed} ditetapkan sebelum eksperimen, diturunkan dengan \textit{SeedSequence} & \textcite{skowronski2026adaptive} \\
\hline
\end{tabular}
\end{table}

Pelatihan dinyatakan konvergen jika tiga syarat terpenuhi bersamaan selama 5.000 langkah, mengikuti \textcite{skowronski2026adaptive}. Pertama, rata-rata bergerak \textit{reward} 10.000 langkah berada dalam rentang $\pm 2\%$ dari rata-rata 20.000 langkah. Kedua, norma gradien di bawah 0,01. Ketiga, kinerja pada episode validasi tidak lagi naik lebih dari 0,01 selama 10.000 langkah. Pada pelatihan ko-evolusioner, syarat keempat ditambahkan, yaitu \textit{exploitability} terhadap kumpulan penyerang tidak berubah lebih dari 5\% antara dua potret berurutan. Syarat ini perlu karena lawan terus berubah sehingga \textit{reward} yang stabil saja belum cukup menunjukkan konvergensi.

\subsection{Estimasi Biaya}
\label{subsec:biaya}
Seluruh perangkat lunak dan data yang digunakan bersifat gratis. Perangkat utama adalah \textit{MacBook} yang sudah dimiliki, sehingga biaya tambahannya hanya listrik. Kebutuhan komputasi diperkirakan pada Tabel~\ref{tab:biaya} dengan acuan satu sesi pelatihan DRL sekitar 8 jam \parencite{skowronski2026adaptive} dan batas 24 jam per sesi pada KNF-05. Jam pada tabel adalah anggaran rencana dan dikalibrasi ulang dengan hasil benchmark pada gerbang G0.

\begin{table}[htbp]
\centering
\caption{Estimasi kebutuhan komputasi}
\label{tab:biaya}
\small
\begin{tabular}{|p{3.8cm}|p{3.4cm}|p{1.9cm}|p{2.7cm}|}
\hline
\textbf{Kegiatan} & \textbf{Perhitungan} & \textbf{Jam komputasi} & \textbf{Lokasi awal} \\
\hline
Verifikasi, kalibrasi, pembanding statis dan berbasis aturan & Hanya CPU & 0 & \textit{MacBook} \\
\hline
Pembanding DRL satu sisi & 10 \textit{seed} $\times$ 8 jam & 80 & \textit{MacBook}, bersyarat G0 \\
\hline
Pelatihan ko-evolusioner utama & 10 \textit{seed} $\times$ 24 jam & 240 & \textit{MacBook}, bersyarat G0 \\
\hline
Penyerang yang dilatih khusus melawan pertahanan beku & 4 pertahanan $\times$ 5 \textit{seed} $\times$ 4 jam & 80 & \textit{MacBook}, bersyarat G0 \\
\hline
Ablasi & 3 varian $\times$ 5 \textit{seed} $\times$ 8 jam & 120 & \textit{MacBook}, bersyarat G0 \\
\hline
\textbf{Total} & & \textbf{520} & \\
\hline
\end{tabular}
\end{table}

Jika dijalankan berurutan di \textit{MacBook}, 520 jam setara sekitar 22 hari tanpa henti. Jendela kalender pelatihan dan eksperimen (minggu 7 sampai 13) mencakup sekitar 49 hari. Dengan asumsi \textit{MacBook} aktif 70\% waktu, kapasitasnya sekitar 34 hari, Jika dijalankan berurutan di \textit{MacBook}, 520 jam setara sekitar 22 hari tanpa henti. Jendela kalender pelatihan dan eksperimen (minggu 7 sampai 13) mencakup sekitar 49 hari. Dengan asumsi \textit{MacBook} aktif 70\% dari waktu kalender, kapasitas efektifnya sekitar 34 hari. Seluruh kegiatan tetap selesai dalam jendela ini selama \textit{MacBook} tidak lebih dari sekitar 1,5 kali lebih lambat daripada anggaran jam. Jika lebih lambat, kegiatan yang melampaui batas tersebut dipindah sesuai gerbang G0. Pelatihan beberapa \textit{seed} secara paralel pada inti yang berbeda dapat menambah margin ini. Batas tersebut diuji pada gerbang G0. Pelatihan beberapa \textit{seed} secara paralel pada inti yang berbeda dapat menambah margin ini. Batas tersebut diuji pada gerbang G0.

Infrastruktur eksternal hanya dipakai untuk kegiatan yang tidak lolos G0. Pilihan pertama adalah GPU kampus atau laboratorium, yang biayanya di luar anggaran penelitian. Jika tidak tersedia, GPU disewa dari penyedia komputasi awan. Dengan tarif sewa RTX 3090 sekitar US\$0,50 per jam, memindahkan seluruh 520 jam membutuhkan sekitar US\$260. Dengan cadangan 30\% untuk pengulangan, batas atas anggarannya sekitar US\$340. Angka ini adalah batas atas, bukan perkiraan utama. Setelah G0, anggaran dihitung ulang sesuai jam yang benar-benar dipindah.
% Sumber tarif: halaman harga RunPod (https://www.runpod.io/gpu-instance/pricing), dicek November 2026. Perbarui sebelum laporan final.

\section{Desain Pengujian dan Evaluasi}
\label{sec:desain_pengujian}
Pengujian dibagi menjadi tiga bagian. Verifikasi memastikan simulator dan agen dibangun dengan benar. Validasi memastikan simulator mewakili Gasper. Eksperimen menjawab RQ2 dan RQ3 serta menguji hipotesis.

\subsection{Verifikasi}
\label{subsec:verifikasi}
Rencana verifikasi disajikan pada Tabel~\ref{tab:verifikasi}. Verifikasi dilakukan sejak tahap implementasi dan diulang setiap kali kode berubah.

\begin{table}[htbp]
\centering
\caption{Rencana verifikasi}
\label{tab:verifikasi}
\small
\begin{tabular}{|p{1.0cm}|p{3.1cm}|p{4.6cm}|p{3.5cm}|}
\hline
\textbf{ID} & \textbf{Objek} & \textbf{Metode} & \textbf{Kriteria lulus} \\
\hline
VF-01 & Justifikasi dan finalisasi & Uji unit pada kasus batas, misalnya suara tepat dua pertiga dan kurang satu suara & Seluruh kasus sesuai Persamaan~\eqref{eq:justified} \\
\hline
VF-02 & \textit{Reward} dan penalti & Hitung manual dibandingkan keluaran simulator untuk \textit{reward} dasar, \textit{slashing} 1/32, batas inklusi, dan \textit{inactivity leak} & Selisih nol \\
\hline
VF-03 & \textit{Action masking} & Uji properti pada 10.000 \textit{state} acak, meniru rangkaian uji kasus tepi \textcite{skowronski2026adaptive} & Tidak ada aksi ilegal yang lolos \\
\hline
VF-04 & Determinisme & Dua kali menjalankan simulasi dan evaluasi kebijakan dengan \textit{seed} yang sama di CPU \textit{MacBook} (\textit{float32}) & Keluaran identik (KNF-03) \\
\hline
VF-05 & Reproduksi model acuan & Konfigurasi 13 validator dan latensi 12,6 detik seperti \textcite{laneve2025stochastic} & Kriteria keberhasilan pertama (KNF-01) \\
\hline
\end{tabular}
\end{table}

VF-05 memerlukan perhatian khusus karena model \textcite{laneve2025stochastic} tidak konsisten soal panjang \textit{epoch}. Karena itu, verifikasi dijalankan pada dua varian, yaitu \textit{checkpoint} setiap 32 \textit{slot} dan setiap 64 \textit{slot}. Varian yang paling mendekati angka acuan dipakai untuk verifikasi, dan perbedaannya dilaporkan sebagai temuan. Eksperimen utama tetap memakai 32 \textit{slot} sesuai Subbab~\ref{sec:gasper}. Determinisme pada VF-04 dinilai pada simulator dan evaluasi di CPU. Pelatihan pada MPS atau CUDA dapat berbeda sedikit antarperangkat, sehingga keterulangannya dinilai secara statistik melalui 10 \textit{seed}, bukan identik bit demi bit.

\subsection{Validasi}
\label{subsec:validasi}
Validasi dilakukan dengan dua cara. Pertama, validasi statistik. Distribusi rasio \textit{slot} kosong, tingkat partisipasi \textit{attestation}, dan rasio blok yatim pada simulator dibandingkan dengan data \textit{on-chain} menggunakan uji Kolmogorov-Smirnov. Simulator dinyatakan valid jika $p > 0{,}05$, mengikuti \textcite{skowronski2026adaptive} yang memperoleh $D = 0{,}032$ dan $p = 0{,}78$ pada simulatornya. Kedua, validasi perilaku. Agen penyerang yang dilatih melawan pertahanan statis diharapkan menemukan kembali pola yang sudah dikenal, misalnya penahanan suara yang mirip \textit{balancing}. Pemeriksaan ini meniru \textcite{wushishi2026coadam}, ketika penyerang yang belajar secara mandiri memilih aksi berdampak tinggi. Jika agen penyerang tidak menemukan pola yang sudah terbukti efektif, ruang aksi atau \textit{reward}-nya perlu diperbaiki sebelum eksperimen dilanjutkan.

Validasi ini tidak menjamin hasil simulasi sama dengan jaringan nyata. \textcite{wushishi2026coadam} menunjukkan bahwa kinerja di uji fisik (0,9975) dapat berbeda jauh dari simulasi (0,8363). Karena itu, kesimpulan tugas akhir ini dibatasi pada lingkungan simulasi, dan ketahanan hasil diperiksa melalui uji sensitivitas terhadap latensi dan skala jaringan.

\subsection{Rancangan Eksperimen}
\label{subsec:rancangan_eksperimen}
Variabel eksperimen disajikan pada Tabel~\ref{tab:variabel}. Eksperimen utama memakai seluruh kombinasi jenis pertahanan dan rezim penyerang pada konfigurasi dasar, yaitu 128 validator dengan porsi \textit{stake} penyerang 0,20. Faktor lain diubah satu per satu pada uji sensitivitas agar jumlah eksperimen tetap terjangkau.

\begin{table}[htbp]
\centering
\caption{Variabel eksperimen}
\label{tab:variabel}
\small
\begin{tabular}{|p{2.4cm}|p{3.2cm}|p{6.6cm}|}
\hline
\textbf{Jenis} & \textbf{Variabel} & \textbf{Level} \\
\hline
Bebas (utama) & Jenis pertahanan & Statis, berbasis aturan, DRL satu sisi, MARL ko-evolusioner \\
\hline
Bebas (utama) & Rezim penyerang & Terskenario (\textit{balancing}, penahanan blok, \textit{censorship}), adaptif selama pengujian, dilatih khusus melawan pertahanan beku \parencite{wushishi2026coadam} \\
\hline
Bebas (sensitivitas) & Porsi \textit{stake} penyerang & 0,10; 0,20; 0,30 \\
\hline
Bebas (sensitivitas) & Jumlah validator & 64, 128, 256, dan 1.000 untuk uji skala \\
\hline
Bebas (sensitivitas) & Latensi & Hasil kalibrasi \textit{on-chain} dan 12,6 detik \parencite{laneve2025stochastic} \\
\hline
Terikat & Metrik keamanan dan biaya & Tabel~\ref{tab:metrik} \\
\hline
Kontrol & Invarian protokol & Ambang dua pertiga, panjang \textit{slot} dan \textit{epoch}, aturan pemilihan rantai \\
\hline
\end{tabular}
\end{table}

Penyerang yang tidak pernah dilihat selama pelatihan, sesuai kriteria keberhasilan kedua, didefinisikan sebagai dua kelompok. Kelompok pertama adalah penyerang terskenario yang tidak pernah menjadi lawan saat pelatihan. Kelompok kedua adalah penyerang yang dilatih dari awal dengan \textit{seed} baru khusus untuk melawan setiap pertahanan yang sudah dibekukan. Kelompok kedua merupakan uji terberat karena penyerang mencari kelemahan satu pertahanan tertentu.

Metrik yang diukur disajikan pada Tabel~\ref{tab:metrik}. Selain rata-rata, kinerja pada kondisi terburuk dilaporkan dengan \textit{Conditional Value at Risk} 5\% (CVaR$_{0{,}05}$), yaitu rata-rata 5\% hasil terburuk, mengikuti \textcite{wushishi2026coadam}.

\begin{table}[htbp]
\centering
\caption{Metrik evaluasi dan keterkaitannya dengan pertanyaan penelitian}
\label{tab:metrik}
\small
\begin{tabular}{|p{3.4cm}|p{6.6cm}|p{2.2cm}|}
\hline
\textbf{Metrik} & \textbf{Definisi} & \textbf{RQ} \\
\hline
Pelanggaran invarian & Jumlah aksi yang melanggar Persamaan~\eqref{eq:batas_parameter}, harus nol & RQ1 \\
\hline
Profit relatif penyerang & Persamaan~\eqref{eq:profit_relatif}, rata-rata dan CVaR$_{0{,}05}$ & RQ2 \\
\hline
Pelanggaran \textit{safety} & Proporsi episode dengan dua \textit{checkpoint} bertentangan yang sama-sama final & RQ2 \\
\hline
Pelanggaran \textit{liveness} & Proporsi \textit{epoch} dengan keterlambatan finalisasi di atas tiga \textit{epoch} & RQ2 \\
\hline
\textit{Exploitability} & Persamaan~\eqref{eq:exploitability} & RQ2, RQ3 \\
\hline
Biaya pertahanan & Penurunan relatif \textit{reward} validator jujur dibanding parameter bawaan & RQ3 \\
\hline
Pola strategi & Frekuensi aksi penyerang, jejak perubahan parameter, dan sebaran aksi pertahanan & RQ3 \\
\hline
\end{tabular}
\end{table}

Ablasi dilakukan pada tiga varian, mengikuti cara \textcite{wushishi2026coadam} menguji sumbangan setiap komponen. Varian pertama tanpa kumpulan penyerang ($p_{\text{pool}} = 0$). Varian kedua tanpa suku biaya pada \textit{reward} pertahanan ($w_4 = 0$). Varian ketiga tanpa \textit{action masking}, yang dijalankan untuk menunjukkan berapa banyak pelanggaran yang terjadi jika batas tidak ditegakkan.

\subsection{Analisis Statistik dan Keputusan Hipotesis}
\label{subsec:analisis_statistik}
Setiap kombinasi eksperimen dijalankan 30 kali dengan \textit{seed} berbeda \parencite{skowronski2026adaptive}. Pertahanan MARL dan pertahanan statis dipasangkan berdasarkan \textit{seed} dan penyerang yang sama, sehingga dipakai uji Wilcoxon \textit{signed-rank} satu sisi. Taraf signifikansi $\alpha = 0{,}05$ dikoreksi dengan Bonferroni sesuai jumlah perbandingan, dan besar efek dilaporkan dengan Cohen's $d$ \parencite{wushishi2026coadam}. Aturan keputusannya adalah sebagai berikut.
\begin{enumerate}
	\item $H_0$ ditolak untuk suatu metrik jika nilai $p$ lebih kecil dari taraf signifikansi yang sudah dikoreksi.
	\item Kriteria keberhasilan kedua terpenuhi jika $H_0$ ditolak untuk profit relatif atau tingkat pelanggaran, dan pelanggaran invarian bernilai nol.
	\item Jika $H_0$ tidak ditolak, hasil tetap dilaporkan beserta besar efeknya sebagai bukti bahwa penyetelan adaptif tidak memberi keuntungan pada kondisi tersebut.
\end{enumerate}

RQ3 dijawab dengan dua analisis. Pertama, \textit{trade-off} keamanan dan biaya dipetakan dengan melatih ulang pertahanan pada beberapa nilai bobot $w_4$, lalu menggambar pasangan profit penyerang dan biaya validator jujur sebagai kurva Pareto. Kedua, strategi penyerang dianalisis dari frekuensi aksinya, seperti analisis frekuensi taktik pada \textcite{wushishi2026coadam}, lalu dibandingkan dengan serangan yang sudah dikenal pada \textcite{laneve2025stochastic}.

\section{Analisis Risiko dan Mitigasi}
\label{sec:risiko}
Risiko yang dapat menghambat penelitian disajikan pada Tabel~\ref{tab:risiko}. Kemungkinan dan dampak dinilai dengan tiga tingkat, yaitu rendah (R), sedang (S), dan tinggi (T). Setiap risiko memiliki mitigasi untuk mencegahnya dan rencana cadangan jika risiko tetap terjadi. Empat risiko yang sudah disebut pada Subbab~\ref{subsec:penentuan_solusi} tercakup pada R2 sampai R5.

\begin{table}[htbp]
\centering
\caption{Analisis risiko dan mitigasi}
\label{tab:risiko}
\footnotesize
\begin{tabular}{|p{0.6cm}|p{2.7cm}|p{0.6cm}|p{0.6cm}|p{3.6cm}|p{3.2cm}|}
\hline
\textbf{ID} & \textbf{Risiko} & \textbf{K} & \textbf{D} & \textbf{Mitigasi} & \textbf{Rencana cadangan} \\
\hline
R1 & Simulator tidak lolos verifikasi atau kalibrasi & S & T & Uji unit sejak awal, verifikasi dua varian panjang \textit{epoch} (Subbab~\ref{subsec:verifikasi}) & Gerbang G1 ditahan. Kalibrasi hanya ke data \textit{on-chain} dan selisih dengan model acuan dilaporkan \\
\hline
R2 & Pelatihan ko-evolusioner berputar tanpa konvergen & S & T & Kumpulan penyerang \parencite{wushishi2026coadam}, empat syarat konvergensi, \textit{replay buffer} diperkecil agar transisi usang cepat terganti & Pelatihan bergantian dengan salah satu agen dibekukan, atau beralih ke keluarga \textit{actor-critic} (Subbab~\ref{subsec:penentuan_solusi}) \\
\hline
R3 & \textit{Exploitability} hanya relatif terhadap kumpulan & T & S & Uji melawan penyerang baru yang dilatih khusus untuk setiap pertahanan beku & Dilaporkan sebagai batas klaim, bukan keseimbangan Nash formal \parencite{wushishi2026coadam} \\
\hline
R4 & Hasil simulasi berbeda dengan jaringan nyata & S & S & Kalibrasi \textit{on-chain} dan uji sensitivitas latensi serta skala & Kesimpulan dibatasi pada simulasi. Uji \textit{testnet} menjadi saran lanjutan \\
\hline
R5 & Satu koalisi tidak mewakili banyak penyerang & T & R & Dicatat sejak asumsi pada Subbab~\ref{subsec:markov_game} & Jika waktu cukup, uji tambahan dengan dua penyerang independen \\
\hline
R6 & Agen penyerang gagal menemukan serangan efektif & S & T & \textit{Reward} bertahap dan inisialisasi dari contoh serangan \textit{balancing} \parencite{laneve2025stochastic} & Evaluasi tetap memakai penyerang terskenario. Jawaban RQ3 dibatasi \\
\hline
R7 & Agen pertahanan menekan penyerang dengan membebani validator jujur & S & S & Suku biaya $w_4$ dan ablasi $w_4 = 0$ & Bobot $w_4$ dinaikkan dan kurva Pareto dilaporkan \\
\hline
R8 & Komputasi melebihi jadwal atau anggaran & S & S & Benchmark pada gerbang G0, uji skala kecil dahulu, jumlah \textit{seed} minimal 10 \parencite{wushishi2026coadam} & Kegiatan yang tidak lolos G0 dipindah ke infrastruktur eksternal dalam batas Subbab~\ref{subsec:biaya}. Ablasi dan uji 1.000 validator dikurangi \\
\hline
R9 & Data \textit{on-chain} terbatas atau akses API dibatasi & R & S & Data diambil sejak awal, disimpan lokal, dan memakai dua sumber & Kalibrasi memakai statistik ringkas publik \\
\hline
R10 & Jadwal terlambat & S & S & Gerbang G1 dan G2, pemantauan mingguan dengan pembimbing & Prioritas pada eksperimen utama. Ablasi dan uji skala dipangkas \\
\hline
R11 & \textit{MacBook} panas, tertidur, atau melambat saat pelatihan panjang & S & S & Adaptor daya terpasang, mode tidur dimatikan, \textit{checkpoint} disimpan berkala agar pelatihan dapat dilanjutkan & Pelatihan dipecah per sesi, kegiatan yang terlambat dipindah sesuai G0 \\
\hline
R12 & Operator PyTorch tidak didukung MPS atau hasilnya menyimpang & S & S & Uji asap pada G0 dengan lapisan standar (\textit{dueling}, ReLU) & Pelatihan dijalankan di CPU atau dipindah ke infrastruktur eksternal \\
\hline
R13 & Perangkat tunggal rusak atau data hilang & R & T & Kode di Git, \textit{checkpoint} dan log dicadangkan berkala ke penyimpanan kedua & Lanjut dari cadangan terakhir pada perangkat pinjaman atau infrastruktur eksternal \\
\hline
\multicolumn{6}{|l|}{K = kemungkinan, D = dampak, R = rendah, S = sedang, T = tinggi} \\
\hline
\end{tabular}
\end{table}
Risiko dengan dampak tinggi, yaitu R1, R2, dan R6, ditangani paling awal karena ketiganya menentukan apakah eksperimen utama dapat berjalan. Jika ketiga risiko tersebut sudah terkendali sebelum minggu ke-11, sisa waktu cukup untuk menyelesaikan eksperimen utama sekalipun ablasi dan uji skala harus dipangkas. Risiko R13 ditangani sejak minggu pertama dengan pencadangan rutin, karena kehilangan data akan menghentikan seluruh pekerjaan.